\cotefile{91}
\scanpage{59}
\unsure{comme par|comme un|canonique}
\[
T\longrightarrow f_*^{\mathrm{ét.}}(T)
\]
qui est tel que l'hom correspondant $\underline{\mathrm{Hom}}_S(\overbrace{f_*^{\mathrm{ét}}(T)}^{\bar T},S')\longrightarrow\underline{\mathrm{Hom}}_S(T,S')$ est un isom. pour tt $S'$. (\unsure{Préciser|Précisions}\,\dots). Il en résulte que pour tt \deleted{$s\in S,$} \unsure{changemt|changement} de base (p. \unsure{ex.} \deleted{$\underline{\mathrm{Hom}}_{k(s)}(\bar T_s,S'_s)\longrightarrow$} \unsure{de} $S=\operatorname{Spec}(\overline{k(s)})$] la $\unsure{=|\simeq}$ relation reste vraie\,; faisant $S'=S\times\underline I$ ($I$ un fini) on en conclut que $T\to f_*^{\mathrm{ét}}(T)$ est surjectif et \inserted{(si $T/S$ est fini)} \unsure{géométriquement} injectif (regarder \deleted{\unsure{sur}} \inserted{au dessus de} chaque pt $s\in S$).

Inversement, si $T\to\bar T$ est un morphisme surjectif et géom injectif de $T$ (fini \inserted{et plat} sur $S$) sur $\bar T$ rev. étale de $S$, alors on a $\underline{\mathrm{Hom}}_S(\bar T,S')=\mathrm{Hom}_S(T,S')$ pour tout $S'$ non ramifié sur $S$ [\unsure{a fortiori|et fini|et finiment}] et comme la $=$ chose reste vrai par ext. de la base \inserted{$S$}, on conclut \deleted{\unsure{Ho}} $\underline{\mathrm{Hom}}_S(\bar T,S')\xrightarrow{\sim}\underline{\mathrm{Hom}}_S(T,S')$, i.e. \unsure{Term.|Terminé|$T$ \illegible{1}}

% archive sheet number: 58
